\section{Introduction}

The integration of object orientation into Tcl has spanned more than two decades,
characterized by ongoing tension between expressive power and philosophical consistency.
\textbf{Incr Tcl (Itcl)}~\cite{dejong,tclwiki-itcl}, developed in the 1990s, introduced
traditional OO features but created fundamental semantic conflicts with the core language.
The community debated standardization through several Tcl Improvement Proposals (TIPs)~%
\cite{tclwiki-tip257}: TIP~\#6 (rejected---include Itcl in core), TIP~\#50 (bundle Itcl
without integration), and TIP~\#257 (TclOO in Tcl~8.6). TIP~\#257~\cite{fellows2006},
led by Donal Fellows, implemented \textbf{TclOO} as a purpose-built OO system for core
integration. Released with Tcl~8.6 in 2013~\cite{tcl86}, TclOO addressed many
compatibility challenges while retaining reference semantics that require explicit object
destruction.

\textbf{Vanilla Object Orientation (VOO)} represents the next stage in this evolution.
Rather than importing OO paradigms from other languages, VOO composes classes from Tcl's
established patterns---lists, dictionaries, namespaces, and procedures. The fundamental
question motivating this work is: \emph{How can object orientation be integrated without
compromising the characteristics that define Tcl's identity?} This question manifests
through several tensions: reference vs.\ value semantics, framework infrastructure
vs.\ compositional patterns, feature completeness vs.\ philosophical alignment, and syntax
familiarity vs.\ conceptual simplicity. The optimal OO framework for Tcl is one that
appears as a natural extension---amplifying the language's inherent strengths rather than
supplanting them.

This paper makes six contributions:
\begin{enumerate}[leftmargin=*]
  \item A novel architecture composing classes from native data structures with automatic
        memory management.
  \item Type-aware field declarations with zero runtime overhead.
  \item Comprehensive benchmarking on Tcl~8.6.13 and~9.0 demonstrating 7--18$\times$
        faster creation and 4--6$\times$ better memory than TclOO, extended to VOO~C++.
  \item The first cross-version scalability analysis of Tcl OO frameworks.
  \item A seamless migration path from lists $\to$ VOO objects $\to$ C++.
  \item Design principles showing that simplicity and performance are complementary
        objectives.
\end{enumerate}
